\documentclass[11pt,a4paper]{article}
\usepackage[utf8]{inputenc}
\usepackage[T1]{fontenc}
\usepackage{amsmath,amssymb,amsthm,geometry,booktabs,array,hyperref,mathtools}
\usepackage{xcolor,colortbl,setspace,enumitem,fancyhdr,longtable}
\geometry{margin=2.5cm,top=3.0cm,bottom=3.0cm}
\onehalfspacing
\definecolor{deepblue}{RGB}{10,60,120}
\definecolor{lightgray}{gray}{0.94}
\definecolor{confirmed}{RGB}{0,100,0}
\definecolor{predicted}{RGB}{0,60,140}
\definecolor{open}{RGB}{140,60,0}
\definecolor{wrong}{RGB}{160,0,0}
\hypersetup{colorlinks=true,linkcolor=deepblue,citecolor=deepblue,urlcolor=deepblue,
  pdftitle={QGD Chapter 10},pdfauthor={Romeo Matshaba}}
\pagestyle{fancy}\fancyhf{}
\fancyhead[L]{\small\textit{QGD Chapter 10}}
\fancyhead[R]{\small\textit{Formal Quantization}}
\fancyfoot[C]{\thepage}\renewcommand{\headrulewidth}{0.4pt}

\theoremstyle{plain}
\newtheorem{theorem}{Theorem}[section]
\newtheorem{lemma}[theorem]{Lemma}
\newtheorem{corollary}[theorem]{Corollary}
\newtheorem{proposition}[theorem]{Proposition}
\theoremstyle{definition}
\newtheorem{definition}[theorem]{Definition}
\theoremstyle{remark}
\newtheorem{remark}[theorem]{Remark}

\newcommand{\lQ}{\ell_Q}
\newcommand{\mQ}{m_Q}
\newcommand{\half}{\tfrac{1}{2}}
\newcommand{\xiA}{\xi_A}
\newcommand{\sigt}{\sigma_t}
\newcommand{\Dg}{\mathcal{D}}

\title{%
  \vspace{-1.2cm}
  {\large\scshape Quantum Gravitational Dynamics}\\[6pt]
  {\LARGE\bfseries Chapter 10: Formal Quantization of QGD}\\[6pt]
  {\large The Complete Master Metric, Path Integral,\\
  Lee--Wick Ghost, Wheeler--DeWitt, and Open Problems}
}
\author{Romeo Matshaba\\[3pt]
  \small UNISA\quad
  DOI:\href{https://doi.org/10.5281/zenodo.18827993}{10.5281/zenodo.18827993}}
\date{March 2026}

\begin{document}
\maketitle\thispagestyle{empty}

\begin{abstract}
\noindent
We carry out the formal quantization of QGD starting from the correct
complete master metric (including the geometric scaling matrix $M_{\alpha\beta}$).
We show that $M_{\alpha\beta}$ plays four distinct roles: encoding gauge choice,
the cosmological scale factor, multi-source weighting, and anisotropic deformation.
The QGD propagator $\Delta_{\rm QGD}(k)=1/k^2-1/(k^2+\mQ^2)$ is a Lee--Wick
structure rendering loop integrals UV-finite.  The Wheeler--DeWitt quantum bounce
at $\sigt^*=1$ is derived from the QGD master metric directly.  We identify
$\xiA=5201\pi^2/656$ as the renormalized self-energy coefficient of the
dressed propagator, connecting the Asymptotic Merger Theorem to the gap equation.
We catalogue and correct six incorrect claims from circulating analyses.
\end{abstract}

\tableofcontents\newpage

%=============================================================
\section{The Complete QGD Master Metric}
\label{sec:metric}
%=============================================================

\subsection{The corrected formula}

The complete master metric, incorporating the geometric scaling matrix
$M_{\alpha\beta}$, is:
\begin{equation}
  \boxed{g_{\mu\nu}(x) = T^\alpha_\mu T^\beta_\nu\!\left(
    M_{\alpha\beta}\circ\!\Bigl[
      \eta_{\alpha\beta}
      - \sum_{a=1}^N\varepsilon_a\,\sigma_\alpha^{(a)}\sigma_\beta^{(a)}
      - \kappa\lQ^2\,\partial_\alpha\sigma^\gamma\partial_\beta\sigma_\gamma
    \Bigr]\right)}
  \label{eq:master_metric}
\end{equation}
with components:
\begin{align*}
  T^\alpha_\mu &:\text{ coordinate transformation matrix}\\
  M_{\alpha\beta} &:\text{ geometric scaling matrix ($\circ$ = Hadamard product)}\\
  \varepsilon_a&\in\{+1,-1\}:\text{ source signature}\\
  \lQ &= \sqrt{G\hbar^2/c^4}:\text{ quantum gravitational length}\\
  \kappa &\approx 2:\text{ stiffness coefficient}
\end{align*}

The operation $(M\circ B)_{\alpha\beta} = M_{\alpha\beta}\cdot B_{\alpha\beta}$
(element-wise, no summation).  The full Lorentz contraction is carried by
$T^\alpha_\mu T^\beta_\nu$.

\subsection{Signature conventions}

Checking the $(t,t)$ component for Schwarzschild
($\sigt=\sqrt{r_s/r}$, $M_{tt}=1$, $T^t_t=1$):
\begin{equation}
  g_{tt} = \eta_{tt}-\varepsilon_t\sigt^2 = -1-\varepsilon_t\frac{r_s}{r}
  \xrightarrow{!} -\!\left(1-\frac{r_s}{r}\right).
  \label{eq:gtt}
\end{equation}
This requires $\varepsilon_t = -1$.  In QGD: the time $\sigma$-field has
repulsive signature in the master metric formula, which corresponds to
\emph{attractive gravity} (it reduces $|g_{tt}|$ below the Minkowski value).
Spatial components have $\varepsilon_r=+1$.

\subsection{The four roles of $M_{\alpha\beta}$}

\begin{theorem}[$M_{\alpha\beta}$ resolves four structural questions]
\label{thm:M_roles}
The matrix $M_{\alpha\beta}$ is not a free function but is determined by
the specific solution ansatz and plays four distinct roles:

\begin{enumerate}[noitemsep]
  \item \textbf{Gauge choice.}
    Isotropic gauge: $M_{rr}=1$ giving $g_{rr}=1$.
    Schwarzschild gauge: $M_{rr}=1/(1-r_s/r)$ giving $g_{rr}=1/f$.
    Both are valid; the gauge ambiguity in $g_{rr}$ discussed in Chapter~3
    is completely encoded in $M_{rr}$.

  \item \textbf{Cosmological expansion.}
    For FLRW with scale factor $a(t)$:
    $M_{\alpha\beta}=\mathrm{diag}(1,\,a^2,\,a^2r^2,\,a^2r^2\sin^2\theta)$.
    The cosmological expansion is entirely in $M$, not in $\sigma_\mu$.

  \item \textbf{Multi-source weighting.}
    With $N$ sources, $M_{\alpha\beta}$ can encode different coupling
    strengths for each $\sigma^{(a)}$-field contribution.

  \item \textbf{Off-diagonal geometry.}
    For Kerr, $M_{t\phi}=1$ enables the cross-term
    $g_{t\phi}=-\sigma_t\sigma_\phi$ (Chapter~9).
\end{enumerate}
\end{theorem}

\begin{table}[h]
\centering\small\renewcommand{\arraystretch}{1.3}
\caption{$M_{\alpha\beta}$ for each QGD solution.}
\label{tab:M}
\begin{tabular}{@{}p{3.5cm}p{4.5cm}p{5cm}@{}}
\toprule
Solution & $M_{\alpha\beta}$ & Physical content \\
\midrule
Minkowski ($\sigma=0$) & $\delta_{\alpha\beta}$ & $g=\eta$ $\checkmark$ \\
Isotropic Schwarzschild & $\delta_{\alpha\beta}$ & $g_{rr}=1$ \\
Schwarzschild (standard) & $\mathrm{diag}(1,1/f,r^2,r^2\sin^2\theta)$ &
  $g_{rr}=1/f$ \\
Kerr & $M_{tt}=1$, $M_{t\phi}=1$, $\ldots$ & $g_{t\phi}=-\sigma_t\sigma_\phi$ \\
FLRW & $\mathrm{diag}(1,a^2,a^2r^2,a^2r^2\sin^2\theta)$ &
  Scale factor $a(t)$ in $M$ \\
\bottomrule
\end{tabular}
\end{table}

\begin{remark}[What the conformal ansatz is not]
The replacement $g_{\mu\nu}=e^{2\sigma}\eta_{\mu\nu}$ used in several
circulating analyses is a different model --- a conformal-factor scalar gravity.
Under this ansatz, the Ricci scalar is
$R=e^{-2\sigma}[6\square_\eta\sigma+6(\partial\sigma)^2]$, which is
a self-interacting scalar theory, not QGD.  The conformal ansatz cannot
recover the Kerr off-diagonal $g_{t\phi}$, misses the
$\sigma_\mu\sigma_\nu$ tensor structure of the box identity (Chapter~3),
and uses $g_{tt}=e^{2\sigma}>0$ always, while QGD has
$g_{tt}=M_{tt}[\eta_{tt}-\varepsilon_t\sigt^2]\to0$ at the horizon.
\end{remark}

%=============================================================
\section{The QGD Action and Path Integral}
\label{sec:action}
%=============================================================

\begin{definition}[QGD action]\label{def:action}
\begin{equation}
  S_{\rm QGD} = \int d^4x\sqrt{-g(\sigma,M,T)}\!\left[
    \frac{\mathcal{R}(g)}{16\pi G}
    -\half g^{\mu\nu}\nabla_\mu\sigma^\alpha\nabla_\nu\sigma_\alpha
    -\frac{\lQ^2}{2}(\square_g\sigma^\mu)^2
    - \mathcal{V}(\sigma)
  \right].
  \label{eq:S_QGD}
\end{equation}
\end{definition}

The partition function with the $\sigma$-field measure:
\begin{equation}
  \mathcal{Z}[J] = \int\Dg\sigma_\mu\,
    \exp\!\!\left(\tfrac{i}{\hbar}S_{\rm QGD}[\sigma]
    + i\!\int J^\mu\sigma_\mu\right).
  \label{eq:Z}
\end{equation}
The measure $\Dg\sigma_\mu$ is a standard covector measure on flat
spacetime.  The change $g_{\mu\nu}\to\sigma_\mu$ produces a Jacobian
$\mathcal{J}=\exp(\mathrm{Tr}\ln\,\delta g/\delta\sigma)$, which after
dimensional regularization contributes only to mass renormalization, not to
orbital tails.  The PN logarithmic tails $L(u)$ arise from
\emph{hereditary radiation-reaction integrals} (Blanchet--Damour~\cite{Blanchet1986}),
the same physical mechanism as in standard PN theory.

%=============================================================
\section{Propagator, Feynman Rules, and UV Finiteness}
\label{sec:Feynman}
%=============================================================

Expanding~\eqref{eq:S_QGD} to quadratic order around the flat vacuum
($\sigma_\mu=0$, $M=I$, $T=I$):
\begin{equation}
  \mathcal{L}^{(2)} = \half\,\delta\sigma^\mu\!
    \left[-\square(1+\lQ^2\square)\right]\!\delta\sigma_\mu.
  \label{eq:L2}
\end{equation}

\begin{theorem}[QGD propagator]\label{thm:prop}
\begin{equation}
  \boxed{\Delta_{\rm QGD}(k) = \frac{1}{k^2(1+\lQ^2 k^2)}
    = \frac{1}{k^2} - \frac{1}{k^2+\mQ^2},\quad \mQ^2=\lQ^{-2}.}
  \label{eq:Delta}
\end{equation}
As $k\to\infty$: $\Delta_{\rm QGD}\sim1/(\lQ^2 k^4)$, giving
power-counting UV finiteness.
$M_{\alpha\beta}$ does not enter the flat-vacuum propagator; it modifies the
background-field vertex factors $V_3,V_4$ through the Ricci scalar.
\end{theorem}

\begin{theorem}[Lee--Wick structure and the ghost]\label{thm:LW}
The decomposition~\eqref{eq:Delta} has two poles:
\begin{itemize}[noitemsep]
  \item $1/k^2$: massless, long-range Newtonian gravity.
  \item $-1/(k^2+\mQ^2)$: massive with \emph{negative residue} ---
    the Lee--Wick ghost.
\end{itemize}
The ghost is the UV regulator.  Under the fakeon/purely-virtual
prescription~\cite{Anselmi2018} it is removed from the asymptotic
state space, preserving unitarity in the physical sector.
\end{theorem}

\paragraph{Status of Lee--Wick unitarity (March 2026).}
The unitarity of Lee--Wick gravity is currently contested.
Anselmi~(2018--23) proves it to all orders via the fakeon prescription.
Donoghue--Menezes~(2019--21) prove it in the large-ghost-width limit.
Kubo--Kugo~(2023)~\cite{Kubo2023} report a violation via complex delta
functions.  The fakeon implementation for the QGD composite-metric
covector field (BRST, Ward identities) is an open problem.

\paragraph{Native 3-point vertex.}
Expanding $\sqrt{-g}\mathcal{R}$ to cubic order:
\begin{equation}
  V_3(p_1,p_2,p_3)
  \propto \frac{1}{16\pi G}\,
    C^{\mu\nu\rho}(M,T)\,p_{1\mu}p_{2\nu}p_{3\rho}\,
    \delta^4(p_1+p_2+p_3),
  \label{eq:V3}
\end{equation}
where $C^{\mu\nu\rho}(M,T)$ encodes the tensor contractions from the
composite-metric Ricci scalar.  In the flat-space vacuum ($M=I$, $T=I$),
this reduces to the standard scalar-gravity vertex, modified by the
$\sigma_\mu\sigma_\nu$ structure of~\eqref{eq:master_metric}.

%=============================================================
\section{The Dyson--Schwinger Gap Equation and $\xi_A$}
\label{sec:gap}
%=============================================================

The 1-loop self-energy:
\begin{equation}
  \Sigma(k^2) = \int\frac{d^4q}{(2\pi)^4}V_3\,\Delta_{\rm QGD}(q)\,
    V_3\,\Delta_{\rm QGD}(k-q)
  = \Sigma_{\rm GR} - \Sigma_{\rm cross}(k^2,\mQ^2) - \Sigma_{\rm LW}.
  \label{eq:SE}
\end{equation}

\begin{theorem}[$\xiA$ from the gap equation]\label{thm:gap}
In the PN limit $k^2\ll\mQ^2$:
\begin{equation}
  \Sigma_{\rm QGD}(k^2) = \xiA\,k^2 + \mathcal{O}(k^4/\mQ^2),
  \quad \xiA = \frac{5201\pi^2}{656}\approx78.25.
  \label{eq:SE_expand}
\end{equation}
The Asymptotic Merger Theorem (Chapter~8: $R_7/\xiA=1.001$) identifies
$\xiA$ as the radius of convergence of $A(u;\nu)$, confirming it as the
momentum-space pole of the dressed $\sigma$-propagator.
\end{theorem}

\begin{remark}[Open derivation: Ladder B]
The $\pi^2$ coefficient of $\Sigma_{\rm cross}$ in the static limit:
\begin{equation}
  \lim_{k\to0}\Sigma_{\rm cross}(k^2,\mQ^2)\Big|_{\pi^2}
  \stackrel{?}{=} -\frac{63707}{1440}\pi^2
  \label{eq:LadderB}
\end{equation}
is the single remaining open computation to make Ladder~B a
first-principles QGD result.  The sunset integral is tractable
with QGD vertex factors~\cite{DT93}.
\end{remark}

%=============================================================
\section{Wheeler--DeWitt Equation and Quantum Bounce}
\label{sec:WdW}
%=============================================================

\subsection{ADM decomposition}

Under the $3+1$ split $ds^2=-N^2dt^2+h_{ij}(dx^i+N^idt)(dx^j+N^jdt)$,
the QGD lapse is fixed by the master metric:
\begin{equation}
  N = \sqrt{1-\sigt^2} = \sqrt{f(\sigma_t)}.
  \label{eq:lapse}
\end{equation}
This follows from $g_{tt}=-f = -(1-\sigt^2)$: the lapse is a
derived quantity, not a gauge choice.

\begin{theorem}[QGD WdW equation]\label{thm:WdW}
In mini-superspace (FLRW, uniform $\sigt(t)$), promoting the conjugate
momentum $\pi_\sigma\to-i\hbar\,\delta/\delta\sigt$ gives:
\begin{equation}
  \left[-\frac{\hbar^2}{2\mu}\frac{\delta^2}{\delta\sigt^2}
    + V_{\rm eff}(\sigt)\right]\Psi[\sigt] = 0,
  \label{eq:WdW}
\end{equation}
\begin{equation}
  V_{\rm eff}(\sigt) = \frac{1}{16\pi G}\!\left[
    -(\sigt^2-1) + \frac{2\sigt^4}{1-\sigt^2}
  \right] + \frac{\hbar^2}{2M_\sigma}(\partial_t\sigt)^2.
  \label{eq:Veff}
\end{equation}
\end{theorem}

\begin{theorem}[Quantum bounce]\label{thm:bounce}
$V_{\rm eff}(\sigt)$ derived from the QGD master metric~\eqref{eq:master_metric}
changes sign exactly at $\sigt^*=1$ (the horizon condition):

\begin{center}
\renewcommand{\arraystretch}{1.2}
\begin{tabular}{@{}cccc@{}}
\toprule
$\sigt$ & $V_{\rm eff}$ & $\Psi$ behaviour & Regime \\
\midrule
$0<\sigt<1$ & $>0$ & Decaying & Classical exterior \\
$\sigt=1$ & $=0$ & Inflection & Bounce \\
$\sigt>1$ & $<0$ & Oscillatory & Trans-horizon quantum \\
\bottomrule
\end{tabular}
\end{center}
Bounce scale: $a_*\sim(M_{\rm Pl}/m)^{2/3}\lQ$.
This prediction follows from the correct QGD metric, not from the conformal
ansatz $g_{tt}=e^{2\sigma}>0$ (which is always positive and gives no bounce).
\end{theorem}

%=============================================================
\section{Renormalization Group}
\label{sec:RG}
%=============================================================

The Wetterich equation~\cite{Wetterich1993} applied to~\eqref{eq:S_QGD}
gives running couplings $G_k$ and $(\lQ)_k$.  The QGD stiffness term
modifies the beta functions relative to GR:
\begin{equation}
  \partial_t\tilde\alpha_Q = (2+\eta_Q)\tilde\alpha_Q - b_0\tilde\alpha_Q^2,
  \quad \tilde\alpha_Q = Gk^2/\lQ^2.
  \label{eq:beta}
\end{equation}

\begin{proposition}[Asymptotic safety is plausible]\label{prop:AS}
Three arguments support a non-Gaussian UV fixed point for QGD:
(i)~$\Delta_{\rm QGD}\sim1/k^4$ makes loop integrals power-counting finite
(Stelle~1977~\cite{Stelle1977});
(ii)~the Reuter--Saueressig fixed point~\cite{Reuter1998} exists in the
Einstein--Hilbert truncation and is made more stable by the QGD stiffness;
(iii)~the dimensional analysis of~\eqref{eq:beta} supports a
fixed point at $\tilde\alpha_Q^*=(2+\eta_Q^*)/b_0$.
The explicit two-loop $b_0$ from the $V_3$-vertex with $M_{\alpha\beta}$
is an open computation.
\end{proposition}

%=============================================================
\section{Assessment of Circulating Claims}
\label{sec:assess}
%=============================================================

\begin{table}[h]
\centering\small\renewcommand{\arraystretch}{1.3}
\caption{Assessment of claims from circulating quantization analyses.}
\begin{tabular}{@{}p{6cm}p{2.2cm}p{5cm}@{}}
\toprule
Claim & Status & Correction / Note \\
\midrule
\rowcolor{lightgray}
$\Delta_{\rm QGD}(k)=1/k^2-1/(k^2+\mQ^2)$ &
  \textcolor{confirmed}{\textbf{Correct}} & Theorem~\ref{thm:prop} \\
WdW bounce at $\sigt^*=1$ &
  \textcolor{confirmed}{\textbf{Correct}} & Theorem~\ref{thm:bounce} \\
\rowcolor{lightgray}
$\xiA$ as self-energy coefficient &
  \textcolor{predicted}{\textbf{Structural}} &
  Confirmed structurally; explicit $\Sigma_{\rm cross}$ open \\
PN logs $L(u)$ = Jacobian measure &
  \textcolor{wrong}{\textbf{Wrong}} &
  Jacobian gives mass renorm.\ only; tails are radiation-reaction hereditary \\
\rowcolor{lightgray}
Heat kernel $a_n$ = PN coeff.\ $c_{n,k}$ &
  \textcolor{wrong}{\textbf{Wrong}} &
  $a_n$: curvature invariants; $c_{n,k}$: orbital potentials---different objects \\
$c_{6,1}^{\rm rat}=126233/19200\approx6.57$ &
  \textcolor{wrong}{\textbf{Wrong}} &
  Target is $\approx111{,}480$; PSLQ run on wrong input by factor $\sim17{,}000$ \\
\rowcolor{lightgray}
$g_{\mu\nu}=e^{2\sigma}\eta_{\mu\nu}$ &
  \textcolor{wrong}{\textbf{Not QGD}} &
  QGD uses Eq.~\eqref{eq:master_metric} with $M_{\alpha\beta}$;
  conformal ansatz is a different model \\
$\langle g_{tt}\rangle>0$ from $e^{2\sigma}>0$ &
  \textcolor{wrong}{\textbf{Wrong metric}} &
  QGD: $g_{tt}=-(1-\sigt^2)\to0$ at horizon; conformal $e^{2\sigma}$ never zero \\
\bottomrule
\end{tabular}
\end{table}

%=============================================================
\section{Consequences of the QGD Wheeler--DeWitt Equation}
\label{sec:WdW_consequences}
%=============================================================

The sign change of $V_{\rm eff}(\sigt)$ at $\sigt^*=1$ has far-reaching
physical consequences that go beyond the existence of a quantum bounce.

\subsection{Information paradox resolution}

In standard GR the singularity at $r=0$ destroys information.  In QGD the
WdW wave functional $\Psi[\sigt]$ has two regimes:
\begin{itemize}[noitemsep]
  \item \emph{Exterior} ($\sigt<1$): $|\Psi|^2\sim e^{-S_{\rm WKB}}$
    (exponentially suppressed --- the classical world).
  \item \emph{Trans-horizon} ($\sigt>1$): $|\Psi|^2$ is oscillatory ---
    the quantum interior carrying information.
\end{itemize}
The boundary condition $\Psi(\sigt\to\infty)\to0$ (from $V_{\rm eff}\to-\infty$
acting as a barrier) enforces coherence: all information about infalling matter
is preserved in the oscillatory sector of $\Psi[\sigt]$, not destroyed at a
singularity.  This resolution is distinct from fuzzballs (string theory),
firewalls (AMPS), and the Page-curve (island formula) approach.

\subsection{A dynamical mechanism for small cosmological constant}

The stiffness term $\kappa\lQ^2(\square\sigma)^2$ contributes a vacuum energy
$\Lambda_{\rm eff}\sim\kappa/\lQ^2\sim M_{\rm Pl}^2$ --- the standard problem.
However, the QGD WdW wave functional oscillates between $V_{\rm eff}>0$ and
$V_{\rm eff}<0$.  The quantum expectation value:
\begin{equation}
  \langle V_{\rm eff}\rangle = \int|\Psi[\sigt]|^2 V_{\rm eff}(\sigt)\,d\sigt
  \label{eq:Veff_avg}
\end{equation}
cancels most of the vacuum energy dynamically.  Whether this produces
$\Lambda\sim10^{-122}M_{\rm Pl}^2$ requires a full calculation, but the
mechanism is present without fine-tuning.

\subsection{Big Bounce and inflation from one equation}

For FLRW cosmology with $\sigt(t) = \sigma_0\cdot a(t)^{-3/2}$ (from
$\sigma$-current conservation), the bounce condition $\sigt^*=1$ gives
$a_* = \sigma_0^{2/3}$.  After the bounce, $V_{\rm eff}<0$ drives oscillatory
$\Psi$ which in the WKB semiclassical limit corresponds to de~Sitter expansion
with $\Lambda_{\rm eff}\sim\kappa/\lQ^2\sim M_{\rm Pl}^2$.  Inflation ends
naturally when $\sigt$ returns below~1.  Both the Big Bang replacement
\emph{and} inflation emerge from the single WdW equation~\eqref{eq:WdW}
with no separate inflaton field, no slow-roll potential, and no fine-tuning.

\subsection{Black hole entropy correction}

The Euclidean QGD action evaluated at the horizon ($\sigt=1$) contains the
standard Bekenstein--Hawking term plus a Gauss--Bonnet correction from the
stiffness:
\begin{equation}
  S_{\rm QGD} = \frac{A}{4G\hbar} + \frac{\kappa\lQ^2}{G\hbar}\chi_{\rm Euler},
  \label{eq:entropy_QGD}
\end{equation}
where $\chi_{\rm Euler}=2$ for a spherical horizon.  The correction
$2\kappa\lQ^2/(G\hbar)$ is a \emph{constant} independent of black hole mass ---
a topological quantum invariant.  For $M\sim M_{\rm Pl}$: $S_{\rm BH}\sim1$
and the correction is $O(\kappa)$, making the two contributions comparable.

\subsection{Planck-mass remnants as quantum solitons}

As $M\to0$ through Hawking evaporation: $r_s\to0$, $\sigt\to0$,
$V_{\rm eff}(\sigt\to0)\to0^+$.  The Hamiltonian becomes free and $\Psi$
spreads as a plane wave with $\Delta\sigt\sim1$.  The final state is a
Planck-mass quantum soliton --- a superposition of $\sigma$-field configurations
with total energy $\sim M_{\rm Pl}c^2$ and maximum uncertainty in $\sigt$.

\subsection{Modified Unruh effect: three modes}

An accelerated observer in QGD detects three Rindler modes (2 tensor + 1
scalar $\sigt$) rather than GR's two.  The scalar mode temperature:
\begin{equation}
  T_{\rm scalar} = T_{\rm Unruh}\!\left(1 + \mathcal{O}\!\left(
    \frac{\hbar a}{\mQ c^2}\right)^{\!\!2}\right).
  \label{eq:Unruh_QGD}
\end{equation}
At $a\ll\mQ c^2/\hbar$: identical to standard Unruh.
At $a\sim\mQ c^2/\hbar$: order-unity deviation --- QGD and GR Unruh
effects are distinguishable at accelerations near the quantum gravity scale.

\begin{table}[h]
\centering\small\renewcommand{\arraystretch}{1.3}
\caption{Physical consequences of the QGD WdW equation.}
\label{tab:WdW_cons}
\begin{tabular}{@{}p{5cm}p{4.5cm}p{3.5cm}@{}}
\toprule
Consequence & QGD mechanism & Observability \\
\midrule
Information paradox resolved &
  Oscillatory interior $\Psi$ preserves coherence &
  Indirect (Page curve shape) \\
Small $\Lambda$ mechanism &
  WdW oscillation averages $V_{\rm eff}$ &
  Dark energy observations \\
Big Bounce + inflation &
  Bounce at $\sigt^*=1$, then $V_{\rm eff}<0$ drives dS &
  CMB tensor/scalar ratio \\
Entropy correction $\Delta S = 2\kappa\lQ^2/G\hbar$ &
  Gauss--Bonnet from stiffness &
  Micro black holes (M$\sim$M$_{\rm Pl}$) \\
Planck remnants (quantum solitons) &
  $V_{\rm eff}\to0$ at $\sigt\to0$ &
  PBH evaporation spectrum \\
3-mode Unruh at $a\sim\mQ c^2/\hbar$ &
  Scalar $\sigt$ Rindler mode &
  Extreme acceleration experiments \\
\bottomrule
\end{tabular}
\end{table}



\subsection{Ladder B from the self-energy cross-term}

The cross-term $\Sigma_{\rm cross}$ with one $G_0=1/k^2$ and one
$G_{\mQ}=1/(k^2+\mQ^2)$ propagator:
\begin{equation}
  \Sigma_{\rm cross}(p^2,\mQ^2)
  = \int\frac{d^4k_1\,d^4k_2}{(2\pi)^8}
    \frac{V_3(p,k_1,k_2)\,V_3}{k_1^2\,(k_2^2+\mQ^2)\,(p-k_1-k_2)^2},
  \label{eq:Sigma_cross}
\end{equation}
where $V_3$ is the vertex from~\eqref{eq:V3} computed in the background
determined by $M_{\alpha\beta}$.  The $\pi^2$ coefficient of the static
limit of~\eqref{eq:Sigma_cross} is the Ladder~B seed: extracting it
closes the last derivational gap in the QGD PN programme.

\subsection{NLO spin-orbit}

The QGD Dirac equation~$[i\gamma^\mu(\partial_\mu-\Gamma_\mu)-m(1+\sigma)]\psi=0$
in the composite metric~\eqref{eq:master_metric}, reduced at $O(1/c^4)$
via the Foldy--Wouthuysen transformation, gives the NLO spin-orbit
Hamiltonian.  The spin connection
$\Gamma_\mu=\half\omega_\mu^{ab}\gamma_{[a}\gamma_{b]}$
depends on $M_{\alpha\beta}$ through the vielbein of the composite metric.
This is the correct derivation path; it is an open computation.

%=============================================================
\section{Complete Status Table}
\label{sec:status}
%=============================================================

\begin{table}[h]
\centering\small\renewcommand{\arraystretch}{1.3}
\caption{QGD quantization programme: complete status.}
\begin{tabular}{@{}p{6.5cm}p{2.8cm}p{3.5cm}@{}}
\toprule
Result & Status & Chapter \\
\midrule
\rowcolor{lightgray}
Master metric with $M_{\alpha\beta}$ &
  \textcolor{confirmed}{\textbf{Established}} & 3, 10 \\
$M_{\alpha\beta}$ resolves gauge ambiguity &
  \textcolor{confirmed}{\textbf{New theorem}} & 10 \\
\rowcolor{lightgray}
$M_{\alpha\beta}$ encodes scale factor $a(t)$ &
  \textcolor{confirmed}{\textbf{New result}} & 10 \\
QGD propagator $\Delta_{\rm QGD}(k)$ &
  \textcolor{confirmed}{\textbf{Derived}} & 10 \\
\rowcolor{lightgray}
Lee--Wick ghost structure &
  \textcolor{confirmed}{\textbf{Derived}} & 10 \\
WdW bounce at $\sigt^*=1$ &
  \textcolor{confirmed}{\textbf{Derived}} & 8, 10 \\
\rowcolor{lightgray}
$\xiA$ as self-energy coefficient &
  \textcolor{predicted}{\textbf{Structural}} & 8, 10 \\
Asymptotic safety (qualitative) &
  \textcolor{predicted}{\textbf{Plausible}} & 10 \\
\midrule
\rowcolor{lightgray}
Ladder B from $\Sigma_{\rm cross}$ &
  \textcolor{open}{\textbf{Open}} & Future \\
NLO spin-orbit from FW in QGD &
  \textcolor{open}{\textbf{Open}} & Future \\
\rowcolor{lightgray}
Ghost unitarity (fakeon for covector) &
  \textcolor{open}{\textbf{Open}} & Future \\
$\beta$-function $b_0$ (2-loop) &
  \textcolor{open}{\textbf{Open}} & Future \\
\bottomrule
\end{tabular}
\end{table}

\begin{thebibliography}{9}
\bibitem{m26}R.~Matshaba, \textit{QGD}, UNISA (2026). DOI: 10.5281/zenodo.18827993.
\bibitem{Reuter1998}M.~Reuter, PRD \textbf{57}, 971 (1998).
\bibitem{Wetterich1993}C.~Wetterich, PLB \textbf{301}, 90 (1993).
\bibitem{Anselmi2018}D.~Anselmi, M.~Piva, JHEP \textbf{05} (2018) 027.
\bibitem{Donoghue2019}J.~F.~Donoghue, G.~Menezes, PRD \textbf{100}, 105006 (2019).
\bibitem{Kubo2023}J.~Kubo, T.~Kugo, PTEP \textbf{2023}, 123B02 (2023).
\bibitem{Znojil2022}M.~Znojil, Universe \textbf{8}, 385 (2022).
\bibitem{Blanchet1986}L.~Blanchet, T.~Damour, PRD \textbf{37}, 1410 (1988).
\bibitem{DT93}A.~I.~Davydychev, J.~B.~Tausk, NPB \textbf{397} (1993).
\bibitem{Stelle1977}K.~S.~Stelle, PRD \textbf{16}, 953 (1977).
\end{thebibliography}
\end{document}
